\documentclass[11pt]{article}

\usepackage[margin=1in]{geometry}
\usepackage[T1]{fontenc}
\usepackage[utf8]{inputenc}
\usepackage{lmodern}
\usepackage{microtype}
\usepackage{amsmath,amssymb,amsthm}
\usepackage{booktabs}
\usepackage{array}
\usepackage{enumitem}
\usepackage{listings}
\usepackage[hidelinks]{hyperref}
\usepackage{xurl}

\newtheorem{definition}{Definition}
\newtheorem{theorem}{Theorem}
\newtheorem{remark}{Remark}

\newcommand{\eps}{\varepsilon}
\newcommand{\TV}{\mathrm{TV}}
\newcommand{\KL}{\mathrm{KL}}

\lstset{basicstyle=\footnotesize\ttfamily,breaklines=true,frame=single,columns=fullflexible}

\newcommand{\figbox}[1]{%
  \begin{center}
  \fbox{\parbox{0.93\linewidth}{\vspace{0.5em}\centering\emph{Figure omitted in the source-only archive.}\\[0.25em]\textbf{#1}\vspace{0.5em}}}
  \end{center}
}

\title{W-Congestion-EQ: Verifiable Congestion Budgets\\
Congestion-Specific Verifier Bundles for Anonymous Overlays}
\author{}
\date{Draft v0.46 (2026-03-23) (archive v690)}

\begin{document}
\maketitle

\begin{abstract}
Congestion-EQ budgets are meaningful only if parties other than the service operator can replay the declared witness/accountant packet and verify that the published knob choices actually bind to it.
We therefore present \emph{W-Congestion-EQ}, a compact congestion-specific verifier bundle for anonymous overlays (and Anonymous-DHT lookup services): a canonicalized claim object binding (i) the declared congestion witness family, (ii) the PSC-Q knob manifest (or other mechanism manifest) that is supposed to meet it, (iii) the accountant arithmetic that yields the stated budget, and (iv) the evidence digests needed to replay that arithmetic offline.

The contribution of this paper is deliberately narrow.
This note is the optional replayable checker companion beneath Congestion-EQ and PSC-Q: it specifies the congestion-specific fields and MUST checks that let an auditor answer the question ``does this exact congestion-budget claim replay under the declared witness, knobs, and accountant inputs?''
Optional operator microstudies, generic notarized attempt logging, transparency anchoring, and downstream release lineage are imported from the PSC-Q addendum, Eval~1, and Release~A rather than re-owned here.
\end{abstract}

\paragraph{Scope.}
We define the \emph{congestion-specific verifier bundle} for Congestion-EQ style claims: canonical claim fields, replayable evidence digests, and the offline checker that recomputes the declared budget.
This paper is intentionally the downstream replayable checker companion rather than the congestion-family entrypoint, and it stops before optional operator microstudies, generic notarized logging, runtime conformance monitoring, or cross-release publication packaging.\cite{series:pscqaddendum,series:evalnotary,series:release,series:abom}

\paragraph{Units.} Budgets are published in bits (use $\log_2$ and $2^{\epsilon}$); results stated in nats can be converted by dividing by $\ln 2$ (see Synthesis~8).\cite{series:notation}

\paragraph{Imports.}
We import the congestion witness family from Congestion-EQ and mechanisms from PSC-Q.\cite{series:congeq,series:pscq}
Generic notarized attempt logging and beacon/transparency binding are treated as owned by Eval~1.\cite{series:evalnotary}
Release-lineage comparison and signed publication receipts are treated as owned by Release~A; packaging/versioning and drift/change control are treated as owned by ABOM/OINL and disagreement-gated recertification.\cite{series:release,series:abom,series:recert}

\paragraph{Exports.}
A \emph{W-Congestion-EQ verifier bundle}: a canonicalized congestion-budget core binding the witness family, mechanism manifest, accountant inputs, and evidence digests, together with only the deterministic checker profile and short replay verdict needed to say whether the declared congestion budget replays offline.

\paragraph{Boundary and earliest sufficient stop.}
This paper is complete once an auditor can canonicalize the bundle, validate the witness/knob/accountant bindings, replay the declared budget arithmetic, and obtain a local checker verdict.
That is enough to own the replayable congestion-checker contract itself. The maintained worked bundle currently mirrors the theorem/mechanism/replay route through the owner map and Synthesis~55 bridge, not through a live \nolinkurl{congestion\_replay\_packet} in \nolinkurl{example\_verifier\_report.json}. This paper therefore remains the exact public source for the checker-level congestion contract until a future artifact issue deliberately emits a verifier-report packet.
If a publisher wants optional operator microstudy tails, notarized attempt logs, split-world-detection receipts, or lineaged public release artifacts, those handoffs belong to the PSC-Q addendum, Eval~1, and Release~A rather than to W-Congestion-EQ.\cite{series:workedexample,series:seriesspines,series:pscqaddendum,series:evalnotary,series:release}

\paragraph{Later-terse-paper citation posture.}
When a later short paper only needs to answer \emph{which replayable congestion-checker contract was declared for this packet?}, it may stop at the minimal public W-Congestion-EQ card below plus the local replay drill, and it may use the owner-map / Synthesis~55 bridge only as the route marker.
It should widen only when the live issue becomes optional operator evidence (the PSC-Q addendum), generic notarized attempt logging (Eval~1), runtime conformance (Operational~B), or signed release lineage (ABOM/OINL and Release~A). The worked example and series spine now already name the narrower replay stop as a route, but this archive cut does not emit a verifier-report replay packet.\cite{series:workedexample,series:seriesspines,series:pscqaddendum,series:evalnotary,series:operationalB,series:abom,series:release}

\paragraph{Artifact policy.}
Source-only; schemas, verifiers, and microstudy plots referenced below are available on request.\cite{series:artifactpolicy}

\paragraph{A minimal public W-Congestion-EQ card.}
Table~\ref{tab:minimal-wcongeq-card} is the smallest public packet later terse papers should need when the live issue is simply \emph{which congestion-specific verifier bundle and replay rule were declared?}

\begin{table}[t]
\centering
\small
\begin{tabular}{@{}p{0.31\linewidth}p{0.60\linewidth}@{}}
\toprule
\textbf{Public row} & \textbf{Pinned value}\\
\midrule
Witness / mechanism binding & \texttt{witness = congeq.wait.twostage.v1}; \texttt{mechanism = pscq.calendar.v1}; same declared anonymous-DHT lookup family as the public Congestion-EQ / PSC-Q packet\\
Canonicalization profile & JCS / RFC~8785 canonical bytes before hashing, signature verification, or outer notarization\\
Accountant profile & \texttt{congeq.epochsum.v1} with stage terms $(0.010,0.008)$ and a five-epoch reset block\\
Replayed budget arithmetic & checker recomputes $\eps_{\mathrm{epoch}}=0.010+0.008=0.018$ and $\eps_{1:5}=5\times 0.018=0.090$\\
Owner-map / evidence boundary & \texttt{example\_line\_item\_owner\_map.json} names the congestion-family replay-adjacent route; no \texttt{bundle://worked-example/congestion-eq} support resolver is emitted in this archive cut\\
Checker profile / verdict & deterministic checker \texttt{wcongeq-checker.v1}; accept iff schema, signature, arithmetic, and support-bundle checks all pass\\
Published claim cap & public congestion-budget claim for this reset block is $0.090$ in TV distance under the declared witness/accountant packet\\
Outer-owner routing & Eval~1 may notarize attempts; Operational~B may pin runtime conformance; ABOM/OINL and Release~A own public claim packaging and lineage\\
\bottomrule
\end{tabular}
\caption{A minimal public W-Congestion-EQ card. This is the non-decorative public answer later terse papers can quote directly when the live issue is the replayable congestion-checker contract rather than extra operator evidence, runtime conformance, or release packaging.}
\label{tab:minimal-wcongeq-card}
\end{table}

This card is intentionally narrow: it pins the witness/mechanism binding, the canonicalization/checker profiles, one concrete replay arithmetic, and the exact accept condition, but it does not re-own optional microstudies, runtime receipts, or signed publication wrappers. The maintained worked bundle currently mirrors the same checker contract only as an owner-map / series-spine route, while \nolinkurl{example\_verifier\_report.json} does not emit a \nolinkurl{congestion\_replay\_packet}. This table therefore remains the local minimal card until a future artifact issue deliberately promotes the replay row into the verifier report.

\paragraph{One tiny owner-map route drill.}
Suppose a later terse paper needs exactly one honest sentence about the already-fixed replay contract. It may now say: \emph{``W-Congestion-EQ declares the replayable checker contract that recomputes the two-stage congestion packet under \nolinkurl{wcongeq-checker.v1} and exports the same five-epoch reset-block cap $0.090$; the maintained owner map records this as the replay-adjacent congestion-family route, but no verifier-report replay packet is emitted in this archive cut.''}
That sentence is honest because the replay contract remains carried by this paper's local card while the worked-bundle route records only ownership and adjacency.

\paragraph{One tiny replay drill.}
With stage terms $(0.010,0.008)$ in a five-epoch reset block, the checker recomputes
\[
\eps_{\mathrm{epoch}} = 0.010+0.008 = 0.018,
\qquad
\eps_{1:5} = 5\times 0.018 = 0.090.
\]
If the canonical bytes, signature, arithmetic, and support-bundle checks also validate, the local verdict is ``accept'' for the published $0.090$ TV cap. A later terse paper can therefore quote the card and this replay arithmetic without reopening the PSC-Q addendum, notarization ledger, or release stack.

\paragraph{PSC-Q admissibility check before replay.}
The checker must also bind the PSC-Q mechanism manifest to the same admissible-substitution rows used by the public PSC-Q card: cover is preemptible by real work, cover does not reserve downstream capacity needed by real work, and the declared utilization interval remains $\rho\in[0.7,0.8]$ with $\rho_{\max}=0.8$.  If those rows are missing or drift, W-Congestion-EQ must reject this fixed checker packet rather than replaying the $0.090$ cap under an additive or contention-bearing successor mechanism.  The rejected packet may still be useful evidence for a new review, but it is not the same public replay contract.

\begin{theorem}[Local replay soundness and non-substitution]
Fix a W-Congestion-EQ verifier bundle whose canonical bytes are computed under the declared JCS/RFC~8785 profile and signed under the declared verification key.
Assume the bundle names witness id \texttt{congeq.wait.twostage.v1}, mechanism id \texttt{pscq.calendar.v1}, accountant profile \texttt{congeq.epochsum.v1}, stage terms $(0.010,0.008)$, a five-epoch reset block, the declared support-bundle digest set, and checker profile \texttt{wcongeq-checker.v1}.
If the checker accepts, then the public cap it accepts is exactly the deterministic replay value
\[
  5\cdot(0.010+0.008)=0.090,
\]
and that cap is bound to the same canonical bytes, signature, witness/mechanism/accountant identifiers, PSC-Q admissibility rows, and support-bundle digests.
Any change to one of those fields changes the replay input and must be treated as a successor checker packet rather than as the same W-Congestion-EQ card.
Conversely, a schema, canonicalization, signature, arithmetic, support-digest, or PSC-Q admissibility failure is a local reject.
\end{theorem}
\begin{proof}
The checker is deterministic by construction.
It parses the declared object only after canonicalization, verifies the signature on those canonical bytes, checks that the pinned witness/mechanism/accountant/checker identifiers and support-bundle digests are exactly the declared ones, and recomputes the epoch and reset-block arithmetic from the signed stage terms.
The only accepted cap for this packet is therefore the recomputed value $5(0.010+0.008)=0.090$.
Changing any pinned field changes either the canonical bytes being signed, the digest set being checked, the arithmetic being replayed, or the PSC-Q admissibility predicate, so the old accept verdict cannot be reused as evidence for the changed packet.
This proves only local replay of the declared claim object; it does not prove that the deployed service followed the declared mechanism, that the support bundle is complete, or that unmodeled congestion/side channels are absent.
\end{proof}

\paragraph{Replay truth boundary.}
A W-Congestion-EQ acceptance verdict is therefore a deterministic receipt for the declared bundle, not a field audit of the operator.
It says ``these signed bytes replay to this congestion-budget claim under this checker,'' not ``the deployment actually produced only those observables.''
Runtime conformance, complete operator evidence, and anti-equivocation logging remain the jobs of Operational~B, the PSC-Q addendum, Eval~1, and Release~A.
This boundary is a release guard: if an operator cannot supply the support digests or conformance evidence named by the bundle, the checker result may be quoted only as a local replay result, not as operational anonymity evidence.

\paragraph{Branched family route.}
When a later terse Anonymous-DHT note needs the whole congestion family but still wants one owner per sentence, it may say: \emph{``Congestion-EQ declares the witness/load-envelope/accountant root with an epoch cap $0.018$ and a five-epoch reset-block cap $0.090$. PSC-Q declares the 2 ms renewal-style calendar, substitution-only slack harvesting, and observer-resolution/utilization packet intended to meet that root. If the live issue is the optional support packet, the PSC-Q addendum separately records the utilization-knee caution and the rule that cover is cancelled or deferred whenever real work arrives. W-Congestion-EQ exports the replayable checker bundle that replays the same $0.090$ cap under the declared witness/mechanism/accountant/support-bundle packet.''}
That route is honest because Congestion~1 still owns the witness/accountant root, PSC-Q still owns the mechanism packet that aims to meet it, the PSC-Q addendum owns only the optional support branch beneath that packet, and this paper owns only the replayable checker bundle that fixes the same claim as a verifier-facing object.\cite{series:congeq,series:pscq,series:pscqaddendum}

\paragraph{One tiny same-card / refreshed-wrapper cutover drill.}
Suppose a later terse paper republishes the same replayable congestion-checker result but only refreshes the outer wrapper: a new release note date or a new lineage notice that still points to the same witness id, mechanism packet id, canonicalization profile, accountant profile and stage terms, support-bundle evidence row, checker profile, and published $0.090$ cap. Then the old W-Congestion-EQ card remains the same replay contract, and later terse notes may keep citing that old card or the matching maintained \nolinkurl{congestion\_replay\_packet} while leaving wrapper lineage to ABOM/OINL and Release~A.
But if the successor changes the witness family id, swaps the mechanism packet id, changes the accountant profile or stage terms, changes the support-bundle evidence row, changes the checker profile, or changes the replay verdict row, the old card fails closed: this is a successor congestion-checker card, not a wrapper-only refresh.
A later terse paper may therefore say only: \emph{``The wrapper lineage was refreshed, but the W-Congestion-EQ checker packet stayed fixed at the previously declared witness / mechanism / accountant / support-bundle / checker contract.''} Anything stronger should reopen this note or one of its neighboring owners.

\paragraph{Minimal reopen ladder.}
Later terse papers should therefore stop at W-Congestion-EQ only while the live question is \emph{which replayable congestion-checker contract was declared for this already-fixed congestion packet?}
The first widening owner is then explicit:
\begin{itemize}[leftmargin=*]
\item reopen \textbf{Congestion~1} if the theorem/accountant root or public congestion witness changed,\cite{series:congeq}
\item reopen \textbf{PSC-Q} if the declared mechanism packet or public calendar/substitution rows changed,\cite{series:pscq}
\item reopen the \textbf{PSC-Q addendum} if the live issue is optional supporting evidence or operator caution slices rather than the replayable checker contract,\cite{series:pscqaddendum}
\item reopen \textbf{Eval~1} if the live issue is generic notarized attempt logging, public randomness, or the gap-free attempt prefix,\cite{series:evalnotary}
\item reopen \textbf{Operational~B} if the live issue is runtime conformance to the declared plan or checker contract,\cite{series:operationalB}
\item reopen \textbf{ABOM/OINL} if the live issue is signed claim-object sameness or public claim-surface wording,\cite{series:abom}
\item and reopen \textbf{Release~A} if the live issue is lineage-tagged publication, signed receipts, or fail-closed release gating.\cite{series:release}
\end{itemize}
This keeps this note narrow: it owns the replayable congestion-checker bundle and its checker-level cutover, but not optional operator evidence, runtime receipts, or release-lineage packaging.

\section{Position in the series}
Congestion-EQ defines the \emph{witness} and \emph{budget metric} (TV constraints over congestion observables) \cite{series:congeq}.
PSC-Q defines a compiler-friendly mechanism family that attempts to meet those budgets under declared calendars and substitution padding \cite{series:pscq}.
This paper (W-Congestion-EQ) is the narrow downstream replayable checker companion: it turns ``we run PSC-Q with knob set $\theta$'' into a canonicalized congestion-specific claim object that a third party can replay offline after Congestion-EQ has already fixed the witness/accountant root and PSC-Q has already fixed the mechanism packet.

The handoff after that replay is intentionally external.
Eval~1 owns notarized attempt logging and beacon/transparency binding \cite{series:evalnotary}; Release~A owns signed release receipts and lineage comparison \cite{series:release}.
The PSC-Q addendum remains only an optional operator-facing evidence tail (microstudies, checklist fragments, and failure-mode notes), while ABOM/OINL names schema versions and reader-facing claim surfaces, and disagreement-gated recertification governs what happens when the knob/evidence surface drifts \cite{series:pscqaddendum,series:abom,series:recert}.

\section{Threat model and goals}
\textbf{Participants.}
A provider runs an anonymity overlay and publishes periodic (e.g., daily) congestion-budget claims.
Users/auditors consume those claims.
\textbf{Adversary.}
We consider a provider who may (i) misreport knobs or accountant outputs, (ii) publish evidence selectively, or (iii) equivocate by presenting inconsistent congestion-budget bundles to different audiences.
\textbf{Goal.}
Define a certificate/receipt interface with an offline verifier that:
(1) checks internal arithmetic consistency; (2) binds claims to declared knobs and witness families; and (3) reduces equivocation via transparency receipts.

\textbf{Non-goals.}
A passing certificate is not a proof of real-world anonymity: it does not prevent hidden side channels.
It is a verifiable statement about what was claimed, how it was computed, and whether the claim is consistent with the public contract objects of the series.

\section{Contract objects and what must be bound}
\subsection{Witness binding (Congestion-EQ import)}
Congestion-EQ represents congestion leakage through a witness $W$ (queueing delay, RTT, loss/ECN marks, or derived aggregates), and budgets the TV distance between witness distributions under different destinations/queries \cite{series:congeq}.
A congestion verifier bundle must bind:
\begin{itemize}[leftmargin=1.2em]
\item the witness family identifier (including units and sampling rule);
\item the policy for truncation/windowing (to avoid unbounded-support pathologies);
\item the budget statement (per epoch and composed) including any reset/flush model.
\end{itemize}

\subsection{Knob binding (PSC-Q import)}
PSC-Q implements congestion equalization with public calendars, substitution padding, and (optionally) declared dithers \cite{series:pscq}.
A congestion verifier bundle must bind the deployed knob set $\theta$, including:
\begin{itemize}[leftmargin=1.2em]
\item calendar structure (period, slotting, any tiered calendars);
\item substitution padding policy (cover distribution; feasibility constraints);
\item retry discipline relevant to congestion observables (schedule-only retries, if any).
\end{itemize}

\subsection{Packaging hooks (ABOM/OINL import)}
We treat the verifier bundle as an ABOM-referenced artifact:
\begin{itemize}[leftmargin=1.2em]
\item \textbf{Schema pinning.} The certificate names a schema version and canonicalization profile.
\item \textbf{Digest binding.} The certificate carries content hashes for any evidence objects referenced (e.g., topology snapshot digest, calibration profile digest) even if those objects are not published in the source-only archive.
\end{itemize}
The point is not to publish everything, but to let a verifier say ``I checked the claim under this exact set of inputs''.

\section{Verifier-bundle format (congestion-specific core)}
We model the congestion-specific checker input as a canonicalized JSON object.
Generic receipt/notarization semantics are imported from Eval~1 and Release~A; here we only specify the congestion-specific \emph{fields and checks} that the local checker consumes.\cite{series:evalnotary,series:release}

\subsection{Minimal field set}
A minimal congestion verifier bundle contains:
\begin{itemize}[leftmargin=1.2em]
\item \texttt{id}: deployment identifier and epoch range;
\item \texttt{witness}: witness family id, sampling and truncation parameters;
\item \texttt{mechanism}: PSC-Q knob manifest $\theta$ (calendar + padding);
\item \texttt{accountant}: per-epoch budget terms and composed $\eps$ (or $(\eps,\delta)$ if used);
\item \texttt{evidence\_digests}: hashes for calibration profiles and reference data;
\item \texttt{signatures}: signer id and signature over canonical bytes;
\item \texttt{receipts} (optional): transparency log receipts for the signed bytes.
\end{itemize}

\subsection{Canonicalization}
JSON has many syntactic representations of the same object.
Therefore the bundle must have a unique canonical byte representation (e.g., RFC~8785/JCS) \cite{rfc8785}.
This is a MUST check: a verifier recomputes canonical bytes before hashing, diffing, or handing the bundle to outer signing/notarization layers.

\section{Mechanical verification checklist}
We separate what the local congestion checker can \emph{mechanically} check from what later notarization or release layers may additionally require.

\subsection{MUST checks}
\begin{enumerate}[leftmargin=1.6em]
\item \textbf{Schema and canonicalization.} Validate the JSON object against a pinned schema version; canonicalize using the declared profile (e.g., JCS) \cite{rfc8785}.
\item \textbf{Signature.} Verify the signature over canonical bytes; verify key identity according to the trust model used by the deployment (e.g., pinned keys, or an identity framework).
\item \textbf{Arithmetic.} Recompute the accountant from the declared per-epoch terms and repetition/reset rules; check that the composed bound matches the certificate.
\item \textbf{Binding.} Check that the witness id, knob manifest, and evidence digests are present and well-formed.
\end{enumerate}

\subsection{Imported notarization / anti-equivocation hooks}
If a deployment additionally publishes notarized attempt logs or transparency receipts, the checker should expose stable digests that those outer layers can bind:
\begin{enumerate}[leftmargin=1.6em]
\item \textbf{Bundle digest.} Hash of the canonical verifier bundle bytes.
\item \textbf{Checker digest.} Hash of the replay verdict and checker profile.
\item \textbf{Schema/version ids.} Stable identifiers for the witness family, knob manifest schema, and accountant profile.
\end{enumerate}
The mechanics of append-only logging, inclusion/consistency proofs, gossip, or cross-release publication are imported from Eval~1 and Release~A rather than specified here.\cite{series:evalnotary,series:release}

\section{Handoffs beyond the verifier bundle}
Once the local checker has replayed the congestion budget, later papers may carry that result forward:
\begin{itemize}[leftmargin=1.2em]
\item \textbf{Eval~1.} Owns gap-free logged attempts, public-randomness binding, and notarized anti-grinding evidence for the evaluation workflow.\cite{series:evalnotary}
\item \textbf{ABOM/OINL.} Owns the reader-facing claim surface and version/diff vocabulary used to describe what changed.\cite{series:abom}
\item \textbf{Release~A.} Owns lineaged signed receipts, PVSA packaging, and cross-release comparison.\cite{series:release}
\item \textbf{Certified~C.} Owns disagreement-gated recertification if the witness/knob/evidence surface changes enough to require a new review packet.\cite{series:recert}
\end{itemize}
\section{Discussion and limitations}
\paragraph{What this does not prevent.}
A provider can satisfy verifier-bundle arithmetic while leaking through channels not modeled in the witness family.
The right interpretation is: ``the published budget is consistent with the public contract objects.''
\paragraph{Why this still helps.}
It turns a congestion claim into an object that can be replayed, diffed, and then handed to later notarization/release layers without redoing the congestion arithmetic each time (like SBOM-style supply-chain artifacts, but for anonymity budgets) \cite{torresarias_intoto_sec19,sigstore_acm2022,slsa_v1}.

\section{Takeaways}
\begin{itemize}[leftmargin=*]
\item \textbf{W-Congestion-EQ is verifiability:} it turns congestion budgets into certificates that third parties can check without trusting the compiler.
\item \textbf{Receipts beat narratives:} publish a verifier report plus inclusion/consistency proofs in an append-only log to reduce equivocation and enable monitoring.
\item \textbf{Interfaces are small:} the certificate binds (i) a concrete plan/manifest and (ii) a dual witness that upper-bounds the congestion objective under the Congestion-EQ accountant.
\end{itemize}

\begin{thebibliography}{99}\small
\bibitem{series:congeq}
Anonymous.
\newblock \emph{Congestion-EQ: Total-Variation Budgets for Congestion Witnesses in Anonymous Overlays}.
\newblock Congestion series Paper~1, The-One-True-Archive (2026).

\bibitem{series:pscq}
Anonymous.
\newblock \emph{PSC-Q: Public Service Calendars and Substitution Padding for Congestion-Equalized Anonymous Lookups}.
\newblock Congestion series Paper~2, The-One-True-Archive (2026).

\bibitem{series:pscqaddendum}
Anonymous.
\newblock \emph{PSC-Q (Addendum): Supplementary Evidence and Operator Microstudies}.
\newblock Congestion series Paper~2A, The-One-True-Archive (2026).

\bibitem{series:operationalB}
Anonymous.
\newblock \emph{Auditable Trace Privacy with Abort}.
\newblock Operational series Paper~B, The-One-True-Archive (2026).

\bibitem{series:evalnotary}
Anonymous.
\newblock \emph{Notarized Anonymity Certificates Under Retries: Grinding-Resistant Evaluation via Transparency Logs, Public Randomness, and $\delta$-Spending}.
\newblock Evaluation series Paper~1, The-One-True-Archive (2026).

\bibitem{series:release}
Anonymous.
\newblock \emph{Privacy as a Release Property: Proof-Carrying Anonymity Receipts for Anonymous DHTs}.
\newblock Release \& destination Paper~A, The-One-True-Archive (2026).

\bibitem{series:abom}
Anonymous.
\newblock \emph{ABOM and OINL: An Anonymity Bill of Materials and an Odds-Inflation Nutrition Label}.
\newblock Anonymity series Paper~C, The-One-True-Archive (2026).

\bibitem{series:recert}
Anonymous.
\newblock \emph{Disagreement-Gated Recertification: Change Control for Certified Anonymous Systems}.
\newblock Certified series Paper~C, The-One-True-Archive (2026).

\bibitem{series:artifactpolicy}
Anonymous.
\newblock \emph{Artifact and Disclosure Policy for the One-True-Archive (Source-Only Public Release)}.
\newblock Synthesis~6, The-One-True-Archive (2026).

\bibitem{series:workedexample}
Anonymous.
\newblock \emph{Worked Example: Receipt Interlock for a Privacy-Preserving Response Path}.
\newblock Series synthesis note (Synthesis~17), 2026. (This archive.)

\bibitem{series:seriesspines}
Anonymous.
\newblock \emph{Series Spines: From Mathematics to Review Handoffs}.
\newblock Series synthesis note (Synthesis~55), 2026. (This archive.)

\bibitem{rfc8785}
A.~Rundgren.
\newblock {RFC}~8785: {JSON} Canonicalization Scheme ({JCS}).
\newblock IETF, 2020.

\bibitem{rfc9162}
B.~Laurie, E.~Kasper, E.~Messeri, and R.~Stradling.
\newblock {RFC}~9162: Certificate Transparency Version~2.0.
\newblock IETF, 2021.

\bibitem{torresarias_intoto_sec19}
S.~Torres-Arias, H.~Afzali, T.~K.~Kuppusamy, R.~Curtmola, and J.~Cappos.
\newblock in-toto: Providing farm-to-table guarantees for bits and bytes.
\newblock In \emph{USENIX Security}, 2019.

\bibitem{sigstore_acm2022}
Z.~Newman, J.~Speed Meyers, and S.~Torres-Arias.
\newblock Sigstore: Software signing for everybody.
\newblock In \emph{ACM CCS}, 2022. DOI: 10.1145/3548606.3560596.

\bibitem{slsa_v1}
SLSA contributors.
\newblock \emph{SLSA Specification v1.0}, 2023.

\bibitem{cose_receipts_draft}
IETF COSE Working Group.
\newblock \emph{COSE Receipts (Internet-Draft)}, 2025.

\bibitem{series:notation}
Anonymous.
\newblock Notation and Minimal Model for the One-True-Archive.
\newblock Series synthesis note (Synthesis~8), 2026. (This archive.)

\end{thebibliography}

\end{document}
